%% 由 build/emit_tex.py 生成（生成物，勿手改）；定制请放 user_style.tex / user_meta.yaml
\chapter{习题5 参考答案}\label{ux4e60ux98985-ux53c2ux8003ux7b54ux6848}

\begin{quote}
说明：第5章（留数理论及其应用）主要习题与参考答案（经典教材整理）。
\end{quote}

\section{习题}\label{ux4e60ux9898-4}

\begin{enumerate}
\def\labelenumi{\arabic{enumi}.}
\item
  下列函数有些什么奇点？若是极点，指出其级：

  \begin{enumerate}
  \def\labelenumii{(\arabic{enumii})}
  \tightlist
  \item
    \(f(z)=\dfrac{1}{z^2(z^2+1)}\)； (2) \(f(z)=\dfrac{\sin z}{z^3}\)；
  \item
    \(f(z)=\dfrac{1}{z^3-z^2-z+1}\)； (4) \(f(z)=\dfrac{\ln(z+1)}{z}\)；
  \item
    \(f(z)=\dfrac{1}{z(e^z+1)}\)； (6) \(f(z)=e^{1/(z-1)}\)；
  \item
    \(f(z)=\dfrac{1}{z^2(e^z-1)}\)； (8)
    \(f(z)=\dfrac{1}{1+z^n}\)（\(n\) 正整数）。
  \end{enumerate}
\item
  证明：若 \(z_0\) 是 \(f(z)\) 的 \(m\)（\(m>1\)）级零点，则 \(z_0\) 是
  \(f'(z)\) 的 \(m-1\) 级零点。
\item
  验证 \(z=\dfrac{i\pi}{2}\) 是 \(\cosh z\) 的一级零点。
\item
  求 \(f(z)=\sin z+\sinh z-2z\) 在 \(z=0\) 处零点的阶数。
\item
  若 \(f(z)\) 与 \(g(z)\) 是以 \(z_0\) 为零点且不恒为零的解析函数，证明
  \[
  \lim_{z\to z_0}\frac{f(z)}{g(z)}=\lim_{z\to z_0}\frac{f'(z)}{g'(z)}
  \] （或两端均为 \(\infty\)）。
\item
  求下列积分的值（\(C\) 取正向）：

  \begin{enumerate}
  \def\labelenumii{(\arabic{enumii})}
  \tightlist
  \item
    \(\oint_C\dfrac{z^2-5z+2}{z(z-1)}dz\)，\(C:|z|=2\)；
  \item
    \(\int_C\dfrac{dz}{z^2-2z+2}\)，\(C:|z-1|=\dfrac32\)；
  \item
    \(\int_0^\infty\dfrac{x^2}{x^4+1}dx\)；
  \item
    \(\int_0^\infty\dfrac{\sin x}{x}dx\)；
  \item
    \(\int_{-\infty}^{+\infty}\dfrac{\cos x}{x^2+4x+5}dx\)；
  \item
    \(\int_0^{2\pi}\dfrac{d\theta}{1+\cos^2\theta}\)。
  \end{enumerate}
\end{enumerate}

\section{参考答案}\label{ux53c2ux8003ux7b54ux6848-4}

\begin{enumerate}
\def\labelenumi{\arabic{enumi}.}
\item
  \begin{enumerate}
  \def\labelenumii{(\arabic{enumii})}
  \tightlist
  \item
    \(z=0\) 二级极点，\(z=\pm i\) 一级极点；
  \item
    \(z=0\) 二级极点（\(\sin z/z^3=1/z^2-1/6+\cdots\)）；
  \item
    \(z=1\) 二级极点，\(z=-1\) 一级极点；
  \item
    \(z=0\) 可去奇点（\(\ln(1+z)/z\to1\)）；
  \item
    \(z=0\) 一级极点；\(z=(2k+1)\pi i\) 一级极点（\(k\neq0\)）；
  \item
    \(z=1\) 本性奇点；
  \item
    \(z=0\) 三级极点（\(e^z-1=z(1+z/2+\cdots)\)，故 \(z^2(e^z-1)\) 是
    \(z\)
    的三级零点，倒数三级极点），\(z=2k\pi i\)（\(k\neq0\)）一级极点；
  \item
    \(z\) 为 \(1+z^n=0\) 的 \(n\) 个根，均为一级极点。
  \end{enumerate}
\item
  设 \(f(z)=(z-z_0)^m\varphi(z)\)，\(\varphi(z_0)\neq0\)，则
  \(f'(z)=(z-z_0)^{m-1}[m\varphi(z)+(z-z_0)\varphi'(z)]\)，括号在
  \(z_0\) 不为零，故为 \(m-1\) 级零点。
\item
  \(\cosh z=\cos(iz)\)；\(\cosh\dfrac{i\pi}{2}=\cos\dfrac{\pi}{2}=0\)；\(\sinh\dfrac{i\pi}{2}=i\sin\dfrac{\pi}{2}=i\neq0\)，故一级零点。
\item
  计算导数：\(f(0)=0,\ f'(0)=0,\ f''(0)=0,\ f'''(0)=0,\ f^{(4)}(0)=0\)?
  实际上 \(\sin z+\sinh z-2z\)
  前几项：\(z-\frac{z^3}{6}+z+\frac{z^3}{6}-2z=0+\frac{-z^3}{6}+\frac{z^3}{6}=0\)，需继续展开：\(\sin z+\sinh z=2z+\frac{z^5}{60}+\cdots\)（奇数幂），故
  \(f(z)=\frac{z^5}{60}+\cdots\)，从而 \(z=0\) 是五级零点。
\item
  设
  \(f(z)=(z-z_0)\varphi(z)\)，\(g(z)=(z-z_0)\psi(z)\)，\(\varphi,\psi\)
  在 \(z_0\) 解析且非零，则
  \(\dfrac{f}{g}=\dfrac{\varphi}{\psi}\)，\(\dfrac{f'}{g'}\) 也趋
  \(\dfrac{\varphi}{\psi}\)；由洛必达形式可得结论。
\item
  \begin{enumerate}
  \def\labelenumii{(\arabic{enumii})}
  \tightlist
  \item
    \(z=0\) 与 \(z=1\) 在 \(|z|=2\) 内：留数法
    \(=2\pi i\left(\lim_{z\to0}zf(z)+\lim_{z\to1}(z-1)f(z)\right)=2\pi i(2-4)=-4\pi i\)；
  \item
    \(z=1\pm i\) 为极点，\(|z-1|=3/2\)
    内两者都在？\(|i|=1<3/2\)，\(|-i|=1<3/2\)，留数和为 \(0\)，故积分为
    \(0\)；
  \item
    \(\dfrac{x^2}{x^4+1}\) 上半平面极点在
    \(e^{i\pi/4},e^{i3\pi/4}\)，留数法得 \(\dfrac{\pi}{\sqrt2}\)；
  \item
    \(\int_0^\infty\dfrac{\sin x}{x}dx=\dfrac{\pi}{2}\)（利用
    \(\dfrac{e^{iz}}{z}\) 的围道积分）；
  \item
    \(R(z)=\dfrac{1}{z^2+4z+5}\) 上半平面仅 \(z=-2+i\)
    一阶极点，\(\int_{-\infty}^\infty\dfrac{\cos x}{x^2+4x+5}dx=\dfrac{\pi}{e}\cos2\)；
  \item
    令
    \(z=e^{i\theta}\)，\(\cos\theta=\dfrac{z+z^{-1}}{2}\)，化为单位圆内留数，得
    \(\int_0^{2\pi}\dfrac{d\theta}{1+\cos^2\theta}=\sqrt2\,\pi\)。
  \end{enumerate}
\end{enumerate}
